\begin{abstract}

Leyva-Vázquez and Smarandache (2025) demonstrated that neutrosophic
T/I/F evaluation, where Truth, Indeterminacy, and Falsity are
independent dimensions not constrained to sum to 1.0, which reveals
``hyper-truth'' (T+I+F > 1.0) in 35\% of complex epistemic
cases evaluated by LLMs. We extend their work in two directions. First,
we replicate and extend their experiment across five model families from
five vendors (Anthropic, Meta, DeepSeek, Alibaba, Mistral), finding
hyper-truth in 84\% of unconstrained evaluations, which confirms the
phenomenon is cross-vendor under our prompt protocol. Second, and
more significantly, we identify a limitation of scalar T/I/F that their
framework cannot address: models adopting an ``Absorption'' position
(T=0, I=1, F=0) produce identical scalar outputs for fundamentally
different epistemic situations (paradox, ignorance, contingency),
collapsing the very distinctions neutrosophic logic was designed to
preserve. We demonstrate that extending the evaluation to include
\textbf{declared losses} (structured descriptions of what the model
cannot evaluate and why) substantially recovers these distinctions. Models
producing identical scalars for paradox and ignorance produce nearly
disjoint loss vocabularies (Jaccard similarity < 0.10 on loss
description keywords), with domain-specific, severity-rated loss
declarations that differentiate the nature of their uncertainty. This
suggests that scalar T/I/F is a necessary but insufficient
representation of epistemic state, and that tensor-structured output
(scalars + losses) provides a more faithful model of LLM epistemic
capabilities.

\end{abstract}
